\usepackage[margin =1in]{geometry}
\usepackage{amssymb,amsthm,amsmath}
\usepackage{enumerate}
\usepackage{hyperref}
\usepackage{cite}
\usepackage[capitalize]{cleveref}
\usepackage[shortlabels]{enumitem}
\usepackage{color}
\usepackage{cancel}
 \usepackage{tikz}
\usepackage{tikz-3dplot}
\usepackage{pgf}
\usepackage{svg}
\usepackage{pgfplots}
\usepackage{import}
\usepackage{svg}
\usepackage{comment}
\usepackage[utf8]{inputenc} % allow utf-8 input
\usepackage[T1]{fontenc}    % use 8-bit T1 fonts
\usepackage{hyperref}       % hyperlinks
\usepackage{url}            % simple URL typesetting
\usepackage{booktabs}       % professional-quality tables
\usepackage{amsfonts}       % blackboard math symbols
\usepackage{nicefrac}       % compact symbols for 1/2, etc.
\usepackage{microtype}
\usepackage{multirow}
\usepackage{xfrac}
\usepackage{todonotes}
\usepackage{titlesec}
\usepackage{caption}
\usepackage{subcaption}

%algorithm packages
\usepackage{algorithm}
\usepackage{algpseudocode}
\usepackage{array}

%% to supress these in algo
\algtext*{EndFor}
\algtext*{EndIf}
\algtext*{EndWhile}
\algtext*{EndProcedure}
\algrenewtext{For}[1]{\textbf{for} #1}
\algrenewtext{While}[1]{\textbf{while} #1}
\algrenewtext{If}[1]{\textbf{if} #1}

\algnewcommand{\Input}{\item[\textbf{Input:}]}
\algnewcommand{\Initialize}{\item[\textbf{Initialize:}]}

\newcommand{\col}{c} %table Macro


\setlength{\parindent}{2em}

\usepackage{varwidth}

\titleformat{\subsubsection}[runin]
{\normalfont\normalsize\bfseries}{\thesubsubsection}{1em}{}


% graphics for images
\usepackage{graphicx}

% If you uncomment the following line, make sure to uncomment the figures
% in the experiments section and comment the ones used now
%\usetikzlibrary{external}
%\tikzexternalize[prefix=figures/]
%\pgfkeys{
%    /tikz/external/mode=list and make
%}

%\usepackage{url}
\usepackage{appendix}
\numberwithin{equation}{section}
%%%%%%%%%% Math %%%%%%%%%%%%%%%

\newcommand{\tr}{\mathrm{Tr}}
\newcommand{\vect}{\mathbf{vec}}

\newcommand{\norm}[1]{\left \| #1 \right \|}
\newcommand{\opnorm}[1]{\norm{#1}_{\mathrm{op}}}
\newcommand{\Rank}{\text{Rank }}
\newcommand{\Ker}{\text{Ker }}
\newcommand{\Ran}{\text{Ran }}
\newcommand{\Dim}{\text{dim }}
\newcommand{\spann}[1]{\mathrm{span}\left\{{#1}\right\}}
\newcommand{\im}{\text{Im }}
\newcommand{\gph}{{\rm gph}\,}
\newcommand{\diag}{{\rm diag}}
\newcommand{\sign}{\mathrm{sign}}



%%%%%%%%% Convex Analysis %%%%%%%%%
\newcommand{\prox}{\text{prox}}
\newcommand{\proj}{\mathrm{proj}}
\newcommand{\dom}[1]{\mathrm{dom }(#1)}
\newcommand{\epi}{\text{epi }}
\newcommand{\hypo}{\text{hypo }}
\newcommand{\interior}{\text{int }}
\newcommand{\bdry}{\text{bdy }}
\newcommand{\relint}{\text{ri }}
\newcommand{\argmin}{\operatornamewithlimits{argmin}}
\newcommand{\mini}{\operatornamewithlimits{minimize}}
\newcommand{\ls}{\operatornamewithlimits{limsup}}
\newcommand{\argmax}{\operatornamewithlimits{argmax}}
\newcommand{\Rext}{(-\infty, \infty]}


%%%%%%%% Theorems %%%%%%%%%%%%%

\theoremstyle{thmstyletwo}%
\newtheorem{thm}{Theorem}[section]
\newtheorem{theorem}[thm]{Theorem}
\newtheorem{definition}[thm]{Definition}
\newtheorem{proposition}[thm]{Proposition}
\newtheorem{lem}[thm]{Lemma}
\newtheorem{lemma}[thm]{Lemma}
\newtheorem{cor}[thm]{Corollary}
\newtheorem{corollary}[thm]{Corollary}
\newtheorem{pa}{Part}
\newtheorem{subclaim}{Subclaim}
{\theoremstyle{plain}
\newtheorem{assumption}{Assumption}}
\newtheorem{model}{Model}


\newtheorem{remark}[thm]{Remark}
\newtheorem{conjecture}[thm]{Conjecture}
\newtheorem{notation}{Notation}
\newtheorem{example}{Example}[section]

\newtheorem{claim}{Claim}
\numberwithin{equation}{section}
\crefname{claim}{claim}{claims}
\Crefname{claim}{Claim}{Claims}
\crefname{lem}{lemma}{lemmas}
\Crefname{lem}{Lemma}{Lemmas}
\crefname{algorithm}{algorithm}{algorithms}
\Crefname{algorithm}{Algorithm}{Algorithms}


% Macros for initialization section
\newcommand{\eps}{\varepsilon}


% Aaron's commands
\newcommand{\conv}[1]{\textrm{conv}\left\{{#1}\right\}}
\newcommand{\C}{\mathbb{C}}
\newcommand{\R}{\mathbb{R}}
\newcommand{\N}{\mathbb{N}}
\newcommand{\Lcal}{\mathcal{L}}
\newcommand{\Fcal}{\mathcal{F}}
\newcommand{\Z}{\mathcal{Z}}
\newcommand{\X}{\mathcal{X}}
\newcommand{\G}{\mathcal{G}}
\newcommand{\Ical}{\mathcal{I}}
\newcommand{\Holder}{H\"older }
\newcommand{\inner}[2]{\left\langle {#1}, {#2} \right\rangle}
\newcommand{\Wopt}{\mathcal{W}_\star^{M,D,N}}
\newcommand{\Lleft}{L^{\leftarrow}}

\newcommand{\ben}[1]{{\color{blue}(#1 -- Ben)}} % Ben's comments
\newcommand{\aaron}[1]{{\color{purple}(#1 -- Aaron)}} % Aaron's comments

\usepackage{mathtools}

\newcommand{\half}[1]{\scalebox{0.5}{#1}}
\usepackage[T1]{fontenc}
\usepackage{lmodern}
\raggedbottom


% PEP macros ---------------
\newcommand{\pAlg}{p^\mathtt{GFOM}}
\newcommand{\dAlg}{d^\mathtt{GFOM}}
\newcommand{\pPEP}[1]{p_{#1}}
\newcommand{\dPEP}[1]{d_{#1}}
\newcommand{\Algclass}[1]{\mathtt{A}_{\mathrm{{#1}}}}
\newcommand{\Pclassd}{\mathtt{P}_{M,D}^d}
\newcommand{\Pclass}{\mathtt{P}_{M,D}}




\newlength{\solutionindent}
\setlength{\solutionindent}{0.1cm}

\newlength{\questionboxrule}
\setlength{\questionboxrule}{0.8pt} % thinner border (default fboxrule is 0.4pt, go lower for thinner)

\newcommand{\solution}[1]{%
    \vspace{0.2cm}
    \noindent\underline{\textit{Solution:}}\\[2pt]
    \begin{adjustwidth}{\solutionindent}{\solutionindent}
        {\color{red} #1}
    \end{adjustwidth}
    \vspace{0.2cm}
}

\newcommand{\algobox}[1]{%
    \begin{center}
        \begin{tikzpicture}
            \node[
                draw=black,
                line width=\questionboxrule,
                rounded corners=1pt,        % adjust for more/less rounding
                fill=white,
                inner sep=8pt,              % controls padding inside the box
                text width=\dimexpr\linewidth-4\fboxsep-2\solutionindent\relax,
            ] {\color{black} #1};
        \end{tikzpicture}
    \end{center}
}

\usepackage[dvipsnames]{xcolor}



\usepackage{hyperref}

\hypersetup{
  pdftitle={
    A Complete Characterization of Optimal 
    Subgradient Methods for Lipschitz Convex Minimization
  },
  pdfauthor={
    Aaron Zoll and Benjamin Grimmer
  },
  pdfsubject={
    A characterization of all fixed-step first-order methods attaining
    the minimax-optimal objective-gap guarantee for Lipschitz convex
    minimization
  },
  pdfkeywords={
    nonsmooth convex optimization,
    Lipschitz convex minimization,
    subgradient methods,
    fixed-step first-order methods,
    minimax optimality,
    optimal first-order methods,
    performance estimation problems,
    worst-case analysis,
    polyhedral characterization,
    first-order oracle complexity
  },
  pdfcreator={LaTeX with hyperref},
  pdfdisplaydoctitle=true,
  unicode=true
}